\section{Abstract}
\label{sec:abstract}

Argus is a harness for running a capable coding-and-reasoning model as an
autonomous research agent over long horizons --- hours to days --- against real,
public benchmarks. Its central claim is not that autonomy \emph{succeeds}, but
that autonomy can be made \emph{reliable and auditable}: a mission either makes
verifiable forward progress under a bounded budget, is explicitly surfaced as
blocked, or is judged complete by a reviewer against a structured checklist ---
never left to spin indefinitely, and never silently ``finished'' without
evidence.

The system rests on one design commitment: \textbf{research judgment belongs to
the agent; the harness is a domain-agnostic pipe}. Four persistent roles ---
Manager, Planner, Engineer, and Reviewer --- carry every scientific decision
(what to try, whether a baseline is strong enough, whether an experiment is
finished, whether to submit). The harness underneath them does only
domain-agnostic work: budget and rate limiting, persistence and working memory,
scheduling and daemon lifecycle, structured input/output, and anti-cheat
guardrails. The harness never uses keyword or regex heuristics to second-guess
the agent about the substance of the research.

This report documents that commitment concretely. It describes the four-role
architecture, in which the Reviewer is the \emph{sole} authority on completion;
the reliability mechanisms that keep a long-running session from degrading
(a reviewer-curated working-memory checkpoint, bounded session rolls, structured
decision-progress classification with a safe round-boundary budget, and
supervised background waits); and an evidence methodology that separates
reproducible in-repository artifacts from operator-observed preliminary
observations, external baselines, and ongoing programs with no public claim.

We are deliberate about what this release does \emph{not} assert. Argus does not
claim clean, publishable Argus scores on the benchmark-reproduction programs it
is actively running (KernelBench, nanoGPT speedrun, nanochat), and it does not
claim to have produced an accepted paper through its autonomous pipeline. External
reference numbers appear only as named third-party targets for comparison, never
as Argus's own measured results. The contribution here is an engineering
substrate for reliable long-horizon autonomy, presented honestly at
\textbf{Technical Report 0.1} maturity.
